\clearpage\pagestyle{otMain}\setlength{\otBlocEspacement}{2.5mm}\fontsize{11}{15}\selectfont\corrigesChapitre{6}{Comparer et calculer des fractions}
\begin{blocLivre}{cor-6-1}\begin{corrige}{1}{Fractions égales}\demarche{Chercher par quel nombre on multiplie ou on divise.}
$\times4$ : $\res{8}$. $\times4$ : $\res{16}$. $:6$ : $\res4$.
\end{corrige}\end{blocLivre}
\begin{blocLivre}{cor-6-2}\begin{corrige}{2}{Simplifier}\demarche{Diviser par un diviseur commun.}
$\frac{15}{20}=\res{\frac34}$ (division par 5) ; $\frac{24}{36}=\res{\frac23}$ (par 12) ; $\frac{45}{10}=\res{\frac92}$ (par 5).
\end{corrige}\end{blocLivre}
\begin{blocLivre}{cor-6-3}\begin{corrige}{3}{Comparer}\demarche{Même dénominateur, même numérateur ou comparaison à 1.}
$\res{\frac57>\frac37}$ ; $\frac56=\frac{10}{12}$ donc $\res{\frac56>\frac7{12}}$ ; $\res{\frac45>\frac49}$ ; $\res{\frac98>1}$, car $9>8$.
\end{corrige}\end{blocLivre}
\begin{blocLivre}{cor-6-4}\begin{corrige}{4}{Calcul mental}\demarche{Penser en quarts.}
$\res{\frac34}$ ; $\res{\frac34}$ ; $\res{\frac14}$ ; $\res1$.
\end{corrige}\end{blocLivre}
\begin{blocLivre}{cor-6-5}\begin{corrige}{5}{Fraction d’un nombre}\demarche{Diviser par le dénominateur, multiplier par le numérateur.}
$12:3\times2=\res8$ ; $10:4\times3=\res{7{,}5}$ ; $35:5=\res7$.
\end{corrige}\end{blocLivre}
\begin{blocLivre}{cor-6-6}\begin{corrige}{6}{Ranger}\demarche{Tout écrire avec le même dénominateur.}
En huitièmes : $\frac{14}8$ ; $\frac{12}8$ ; $\frac98$ ; $\frac{15}8$. Ordre : $\res{1+\frac18<\frac32<\frac74<\frac{15}8}$.
\end{corrige}\end{blocLivre}
\begin{blocLivre}{cor-6-7}\begin{corrige}{7}{Sommes et différences}\demarche{Écrire d’abord avec un même dénominateur.}
$\frac68+\frac58=\res{\frac{11}8}$ ; $\frac{11}{12}-\frac4{12}=\res{\frac7{12}}$ ; $\frac8{12}+\frac3{12}=\res{\frac{11}{12}}$.
\end{corrige}\end{blocLivre}
\begin{blocLivre}{cor-6-8}\begin{corrige}{8}{Multiplier par un entier}\demarche{Multiplier le numérateur par l’entier.}
$\res{\frac74}$ ; $\res{\frac{10}3}$ ; $\frac{30}6=\res5$ ; $\res{\frac74}$ : le quart de 7 est égal à 7 fois un quart.
\end{corrige}\end{blocLivre}
\begin{blocLivre}{cor-6-9}\begin{corrige}{9}{Sur une même droite}\demarche{Compter les pas, puis simplifier.}
A : $\frac9{12}=\res{\frac34}$. B : $\frac{10}{12}=\res{\frac56}$. C : $\frac{16}{12}=\res{\frac43}$. Ordre : $\res{\frac34<\frac56<\frac43}$.
\label{fig-cor-6-9}\figureGeo[corrige]{fc-9}
\end{corrige}\end{blocLivre}
\begin{blocLivre}{cor-6-10}\begin{corrige}{10}{Égales ou pas ?}\demarche{Simplifier les deux fractions.}
$\frac69=\frac23$ et $\frac{10}{15}=\frac23$ : $\res{\text{égales}}$. $\frac46=\frac23=\frac8{12}$ et $\frac68=\frac34=\frac9{12}$ : $\res{\text{différentes}}$.
\end{corrige}\end{blocLivre}
\begin{blocLivre}{cor-6-11}\begin{corrige}{11}{Le gâteau}\demarche{Additionner en douzièmes.}
$\frac14+\frac13=\frac3{12}+\frac4{12}=\res{\frac7{12}}$. Il reste $1-\frac7{12}=\res{\frac5{12}}$ du gâteau.
\end{corrige}\end{blocLivre}
\begin{blocLivre}{cor-6-12}\begin{corrige}{12}{La cantine}\demarche{Appliquer la fraction au nombre d’élèves.}
$\frac34\times28=28:4\times3=\res{21}$ élèves. Fraction restante : $1-\frac34=\res{\frac14}$, soit $\res7$ élèves.
\end{corrige}\end{blocLivre}
\begin{blocLivre}{cor-6-13}\begin{corrige}{13}{La recette}\demarche{Multiplier une fraction par un entier.}
12 personnes, c’est 3 fois plus : $3\times\frac34=\res{\frac94}$ L, soit $\res{2{,}25}$ L.
\end{corrige}\end{blocLivre}
\begin{blocLivre}{cor-6-14}\begin{corrige}{14}{Le réservoir}\demarche{Calculer en litres, puis contrôler avec les fractions.}
Au départ : $\frac56\times60=50$ L. On utilise $\frac14\times60=15$ L : il reste $\res{35}$ L. Fraction : $\frac56-\frac14=\frac{10}{12}-\frac3{12}=\res{\frac7{12}}$, et $\frac7{12}\times60=35$.
\end{corrige}\end{blocLivre}
\begin{blocLivre}{cor-6-15}\begin{corrige}{15}{Inventer un problème}\demarche{Partir du calcul pour imaginer une situation.}
Exemple : « Dans un club de 30 enfants, les deux cinquièmes font du judo. Combien d’enfants font du judo ? » Réponse : $30:5\times2=\res{12}$ enfants. Toute situation où l’on prend les $\frac25$ d’une quantité de 30 convient.
\end{corrige}\end{blocLivre}
